\chapter{Grafy Taita, wielomian Tutte'a i macierze Goeritza}
\label{ch:tg-tait-goeritz}

W rozdziale~\ref{ch:tg-tutte} i w sekcji~\ref{sec:kt-jones}
poznaliśmy dwa opisy diagramu węzła: rozkład płaszczyzny
na obszary oraz wybory wygładzeń w nawiasie Kauffmana.
Zaskakująco wiele informacji o diagramie można przechować
w grafie planarnym: obszary stają się wierzchołkami,
a skrzyżowania krawędziami. Ten rozdział pokazuje
jednocześnie, kiedy opis grafowy jest pełny, a kiedy
trzeba zachować także znaki lub dane zanurzenia.

\section{Kolorowanie szachownicowe i grafy Taita}
\label{sec:tg-tait-def}

\begin{definicja}[Graf Taita]
Niech $D$ będzie spójnym regularnym diagramem splotu
na orientowanej sferze $S^2$, z co najmniej jednym
skrzyżowaniem. Kolorujemy składowe $S^2\setminus D$
na czarno i biało tak, aby obszary sąsiadujące przez
odcinek diagramu miały różne kolory. \emph{Czarny graf
Taita} $G_B(D)$ ma jeden wierzchołek w każdym czarnym
obszarze oraz jedną krawędź za każde skrzyżowanie,
łączącą wierzchołki dwóch czarnych sektorów przy tym
skrzyżowaniu. Każdą krawędź prowadzimy od wierzchołka
przez wnętrze czarnego obszaru, następnie przez małe
otoczenie skrzyżowania i przez drugi czarny obszar
do odpowiedniego wierzchołka. Łuki wewnątrz obszarów
dobieramy bez przecięć poza wspólnymi wierzchołkami.
Analogicznie definiujemy biały graf Taita $G_W(D)$.
Pętle i wielokrotne krawędzie są dozwolone.
\end{definicja}

\begin{stwierdzenie}[Dualność grafów Taita]
Dla każdego spójnego regularnego diagramu splotu $D$
z co najmniej jednym skrzyżowaniem istnieje kolorowanie
szachownicowe jego obszarów, jednoznaczne z dokładnością
do zamiany kolorów. Po ustaleniu kolorowania
grafy $G_B(D)$ i $G_W(D)$ są planarnymi grafami dualnymi.
\end{stwierdzenie}
\begin{proof}
Na sferze diagram jest zanurzonym grafem czteroregularnym.
Każdy jego wierzchołek ma parzysty stopień, toteż
twierdzenie \ref{thm:tg-face-two-color} daje kolorowanie
szachownicowe. Istotna jest parzystość stopni
wierzchołków diagramu, nie długości obwodów jego ścian.
Graf sąsiedztwa obszarów przez łuki diagramu jest spójny:
jest grafem dualnym spójnego grafu diagramu na sferze
(por. stwierdzenie \ref{thm:tg-dual-basic}).
Po ustaleniu koloru jednego obszaru kolor każdego
innego jest więc wymuszony przez przechodzenie wzdłuż
dowolnej ścieżki w tym grafie. Istnienie kolorowania
zapewnia zgodność tych wymagań dla różnych ścieżek.
Są dokładnie dwa kolorowania, różniące się zamianą
nazw kolorów.

W każdym skrzyżowaniu czarne sektory stanowią przeciwległą
parę, a białe drugą. Wybieramy w każdym czarnym obszarze
punkt i prowadzimy od niego rozłączne (poza tym punktem)
łuki do wystąpień czarnych sektorów przy skrzyżowaniach.
Jest to możliwe w modelu dysku charakterystycznego ściany
spójnego grafu na sferze (por. dowód twierdzenia
\ref{thm:tg-euler-plane}). Łącząc łuki parami przez
skrzyżowania, otrzymujemy krawędzie czarnego grafu.
Analogicznie konstruujemy biały graf, którego krawędzie
przecinają krawędzie czarne dokładnie w odpowiadających
im skrzyżowaniach.

Czarny graf jest spójny: dwa kolejne skrzyżowania
wzdłuż łuku diagramu mają po tej samej stronie łuku
wspólny czarny obszar, więc odpowiadające im krawędzie
czarnego grafu mają wspólny koniec. Spójność diagramu
łączy w ten sposób dowolne dwie krawędzie czarnego
grafu. Każdy czarny obszar jest incydentny z pewnym
skrzyżowaniem, gdyż diagram jest spójny i ma co najmniej
jedno skrzyżowanie. Nie ma więc izolowanych wierzchołków
czarnego grafu. Ten sam argument działa dla białego.

Regularne otoczenie czarnego grafu tworzymy z dysków
w czarnych obszarach oraz pasów przy skrzyżowaniach.
Jego dopełnienie składa się z dysków odpowiadających
białym obszarom: lokalnie przy każdym skrzyżowaniu
czarny pas oddziela oba białe sektory, a wzdłuż łuku
diagramu brzeg otoczenia biegnie między czarnym i białym
obszarem. Każdy dysk dopełnienia jest zatem białym dyskiem
charakterystycznym z kołnierzem wzdłuż jego brzegu.
Po przejściu od otoczenia do samego czarnego grafu
otrzymujemy ściany odpowiadające białym obszarom.
Wobec konstrukcji krawędzi jest to dualność komórkowa
z rozdziału \ref{ch:tg-powierzchnie}.
\end{proof}

\begin{przyklad}[Dwa grafy Taita trójlistnika]
Dla standardowego trójskrzyżowaniowego diagramu trójlistnika
jeden wybór kolorowania daje graf trójkątny $C_3$.
Drugi daje graf o dwóch wierzchołkach połączonych trzema
równoległymi krawędziami. Są to wzajemnie dualne grafy
planarne; w obu przypadkach liczba krawędzi wynosi trzy.
Sama liczba krawędzi nie wystarcza do odtworzenia splotu:
w ogólności niezbędne są znaki skrzyżowań i sposób
zanurzenia grafu.
\end{przyklad}

\begin{center}
\begin{tikzpicture}[scale=.9,line cap=round,>=Stealth]
 \begin{scope}[shift={(-2.65,0)}]
  \node[tgBlueVertex] (a) at (0,1.1) {};
  \node[tgBlueVertex] (b) at (-1,-.65) {};
  \node[tgBlueVertex] (c) at (1,-.65) {};
  \draw[tgAccent] (a)--(b)--(c)--cycle;
  \node[font=\small,text=tgNavy] at (0,-1.25) {czarny graf Taita: $C_3$};
 \end{scope}
 \draw[tgTeal,-{Stealth[length=3mm]},line width=1.4pt] (-.85,.2)--(.85,.2);
 \begin{scope}[shift={(2.65,0)}]
  \node[tgRedVertex] (u) at (-.7,.05) {};
  \node[tgRedVertex] (v) at (.7,.05) {};
  \draw[tgTree] (u)--(v);
  \draw[tgTree] (u) to[bend left=55] (v);
  \draw[tgTree] (u) to[bend right=55] (v);
  \node[font=\small,text=tgNavy] at (0,-1.25) {biały graf Taita: trzy krawędzie};
 \end{scope}
\end{tikzpicture}
\end{center}

\section{Nawias Kauffmana ze zbiorów krawędzi}
\label{sec:tg-tait-tutte}

\begin{lemat}[Jednolitość wygładzeń w diagramie alternującym]
\label{lem:tg-tait-alternating}
Niech $D$ będzie spójnym alternującym diagramem splotu
z co najmniej jednym skrzyżowaniem. Dokładnie jeden
z dwóch wyborów nazw kolorów w kolorowaniu szachownicowym
ma tę własność, że wygładzenie A przy każdym
skrzyżowaniu łączy czarne sektory.
\end{lemat}
\begin{proof}
Mówimy, że wygładzenie łączy sektory, jeśli po usunięciu
skrzyżowania te dwa przeciwległe sektory należą lokalnie
do tego samego obszaru między wygładzonymi łukami.
W konwencji z sekcji \ref{sec:kt-jones}, gdy pasmo
z lewego górnego do prawego dolnego końca przechodzi
górą, wygładzenie A jest pionowe: łączy sektory górny
i dolny. Sektory te mają ten sam kolor.

Skierujmy jeden łuk diagramu od skrzyżowania $c$ do
skrzyżowania $c'$ (dopuszczamy $c=c'$). Lewą i prawą
stronę łuku rozumiemy względem jego kierunku i orientacji
sfery. Przy przejściu przez skrzyżowanie wzdłuż pasma
górnego sektor łączony przez A leży po lewej stronie
pasma na wejściu, a po prawej na wyjściu. Dla pasma
dolnego jest odwrotnie: po prawej na wejściu, po lewej
na wyjściu. Istotnie, dla pasma górnego skierowanego
z lewego górnego do prawego dolnego końca są to
odpowiednio sektor górny i dolny. Dla pasma dolnego
skierowanego z prawego górnego do lewego dolnego
końca są to również odpowiednio sektor górny i dolny,
ale po przeciwnych stronach skierowanego pasma.
Odwrócenie kierunku pasma zamienia wejście z wyjściem
oraz lewą stronę z prawą; obrót lokalnego rysunku
zachowuje powyższe reguły.

Ponieważ diagram jest alternujący, jeśli łuk wychodzi
z pasma górnego przy $c$, to wchodzi na pasmo dolne
przy $c'$. Sektor A leży wtedy po prawej stronie
łuku przy obu końcach. Jeśli łuk wychodzi z pasma
dolnego, to wchodzi na pasmo górne i sektor A leży
po lewej stronie przy obu końcach. Kolor danej strony
łuku nie zmienia się wzdłuż jego wnętrza. Zatem kolor
sektorów łączonych przez A jest taki sam przy $c$ i $c'$.
Odwrotnie, gdy oba końce łuku są typu górnego albo
oba typu dolnego, sektory A leżą po przeciwnych
stronach łuku i mają różne kolory. W szczególności
jednolitość koloru sektorów A przy wszystkich
skrzyżowaniach wymusza alternowanie.
Spójność grafu diagramu pozwala przenieść tę równość
na wszystkie skrzyżowania. Zamiana nazw kolorów pozwala
uzyskać kolor czarny i tylko jeden z dwóch wyborów nazw
ma tę własność.
\end{proof}

\begin{twierdzenie}[Wielomian Tutte'a a nawias diagramu alternującego]
\label{thm:tg-tait-tutte}
Niech $D$ będzie spójnym, niezorientowanym, alternującym
diagramem splotu z co najmniej jednym skrzyżowaniem,
a $G$ jego grafem Taita dla kolorowania
szachownicowego wybranego tak, by w każdym skrzyżowaniu
wygładzenie A łączyło czarne sektory. Niech
$v=|V(G)|$ oraz $e=|E(G)|$. W przyjętej wcześniej
normalizacji nawiasu Kauffmana, w której
$\langle\bigcirc\rangle=1$ i
$\delta=-A^2-A^{-2}$, zachodzi
\[
\langle D\rangle
=A^{2v-e-2}\,T_G(-A^{-4},-A^4).
\]
Dla diagramu z orientacją wielomian Jonesa otrzymujemy
przez dodatkową normalizację
$V_L(t)=(-A^3)^{-w(D)}\langle D\rangle\big|_{A=t^{-1/4}}$.
\end{twierdzenie}

\begin{proof}
Niech $F\subseteq E(G)$ oznacza krawędzie tych skrzyżowań,
w których wybieramy wygładzenie A; w pozostałych
wybieramy B. Tak określony stan ma współczynnik
$A^{|F|}A^{-(e-|F|)}=A^{2|F|-e}$.
Czarne obszary diagramu można potraktować jako dyski,
a krawędzie $F$ jako wąskie pasma łączące te dyski.
Otrzymana płaska powierzchnia regularnego otoczenia
podgrafu rozpinającego $(V(G),F)$ ma $k(F)$ składowych,
charakterystykę Eulera $v-|F|$ i, ponieważ leży na
sferze, rodzaj zero na każdej składowej. Jeżeli
$b(F)$ oznacza liczbę jej okręgów brzegowych, to
$2k(F)-b(F)=v-|F|$, czyli
$b(F)=2k(F)-v+|F|$.
Poza dyskami skrzyżowań brzeg regularnego otoczenia
biegnie wzdłuż łuków diagramu. W dysku skrzyżowania
należącego do $F$ pas łączy oba czarne sektory,
a jego dwie boczne krawędzie tworzą wygładzenie A.
Jeżeli skrzyżowanie nie należy do $F$, pasu nie ma:
brzegi dysków czarnych sektorów pozostają osobno,
co daje wygładzenie B. Ta lokalna odpowiedniość
skleja się wzdłuż wszystkich łuków diagramu i daje
bijekcję okręgów brzegu z okręgami stanu Kauffmana.
W szczególności $b(F)\geq1$, gdyż otoczenie jest
niepuste i ma niepuste dopełnienie w sferze.
Zatem
\[
\langle D\rangle=
\sum_{F\subseteq E(G)}
 A^{2|F|-e}\delta^{2k(F)-v+|F|-1}.
\]
Dla spójnego grafu $G$ wielomian Tutte'a można
zapisać jako sumę po podzbiorach krawędzi:
\[
T_G(x,y)=\sum_{F\subseteq E(G)}
(x-1)^{k(F)-1}(y-1)^{|F|-v+k(F)}.
\]
Ponieważ przy $x=-A^{-4}$ i $y=-A^4$ zachodzi
$x-1=A^{-2}\delta$ i $y-1=A^2\delta$,
otrzymujemy po podstawieniu
\[
A^{2v-e-2}T_G(-A^{-4},-A^4)
=\sum_{F\subseteq E(G)}
 A^{2|F|-e}\delta^{2k(F)-v+|F|-1},
\]
co jest dokładnie poprzednią sumą stanów.
Istnienie użytego kolorowania wynika z lematu
\ref{lem:tg-tait-alternating}. Gdy zamiast $G$ wybierzemy
graf biały $G^*$, to $v^*=e-v+2$ ze wzoru Eulera,
a z twierdzenia \ref{thm:tg-tutte-dual} mamy
$T_{G^*}(x,y)=T_G(y,x)$. Ten sam nawias zapisuje się więc
jako $A^{-(2v^*-e-2)}T_{G^*}(-A^4,-A^{-4})$.
W przypadku spójnego diagramu bez skrzyżowań mamy
po prostu pojedynczy okrąg i nawias równy $1$.
Formuła na $V_L$ jest definicją normalizacji
podaną w części poświęconej wielomianowi Jonesa.
\end{proof}

\begin{przyklad}[Pełny rachunek dla trójlistnika]
Dla grafu trójkątnego $G=C_3$ wyznaczyliśmy wcześniej
$T_G(x,y)=x^2+x+y$. Dlatego
\[
T_G(-A^{-4},-A^4)=A^{-8}-A^{-4}-A^4,
\quad
\langle D\rangle=A^{2\cdot3-3-2}T_G(-A^{-4},-A^4)
=A^{-7}-A^{-3}-A^5.
\]
Dla dodatnio zorientowanego diagramu trójlistnika
$w(D)=3$, zatem
$(-A^3)^{-3}\langle D\rangle
=A^{-4}+A^{-12}-A^{-16}$.
Po przyjęciu $t=A^{-4}$ dostajemy
$V_T(t)=t+t^3-t^4$, zgodnie z obliczeniem
z algebry Temperleya--Lieba z sekcji~\ref{sec:kt-jones}.
Dwa zupełnie różne
rachunki prowadzą więc do tego samego niezmiennika.
\end{przyklad}

\section{Macierz Goeritza i wyznacznik węzła}
\label{sec:tg-goeritz}

\begin{definicja}[Podpisany laplasjan grafu Taita]
Dla wybranego kolorowania szachownicowego diagramu $D$
przypisujemy każdemu skrzyżowaniu wagę
$\eta(c)\in\{-1,1\}$: znak dodatni oznacza,
że lokalne wygładzenie A łączy czarne sektory,
a ujemny, że łączy białe. W lokalnym rysunku z sekcji
\ref{sec:kt-jones}, gdy pasmo górne biegnie od lewego
górnego do prawego dolnego końca, znak $+1$ odpowiada
czarnym sektorom górnemu i dolnemu, a $-1$ sektorom
lewemu i prawemu. Nie jest to znak zorientowanego skrzyżowania.

Niech $G=G_B(D)$ ma $n$ wierzchołków i $m$ krawędzi.
Ustalamy orientację pomocniczą i macierz incydencji
$B\in M_{n\times m}(\mathbb Z)$, której kolumna dla
krawędzi $u\to v$ jest równa $[v]-[u]$.
Zachowujemy zerową kolumnę każdej pętli i osobną
kolumnę każdej krawędzi równoległej. Macierz wag ma
rozmiar $m\times m$ i tę samą kolejność krawędzi.
Definiujemy macierz
\[
\mathcal G_D=B\,\operatorname{diag}(\eta(c))\,B^T.
\]
Jej elementy pozadiagonalne są ujemnymi sumami wag
krawędzi łączących odpowiednie wierzchołki; przekątna
jest sumą wag incydentnych krawędzi niebędących pętlami.
Pętle nie wnoszą wkładu do żadnego elementu macierzy.
Odwrócenie orientacji pomocniczej zmienia znak kolumny
$b_c$, ale nie składnika $\eta(c)b_cb_c^T$, zatem nie
zmienia macierzy. \emph{Zredukowaną
macierzą Goeritza} nazywamy macierz otrzymaną przez
usunięcie jednego wiersza i odpowiadającej kolumny
z $\mathcal G_D$.
\end{definicja}

\begin{stwierdzenie}[Suma drzew z wagami]
\label{prop:tg-tait-signed-tree}
Niech $G$ będzie skończonym spójnym multigrafem
o wierzchołkach $1,\ldots,n$, $n\geq1$, z dopuszczeniem
pętli i wagami $w_e\in\mathbb R$. Niech $B$ będzie
jego macierzą incydencji z osobną kolumną każdej krawędzi.
Dla macierzy $L=B\operatorname{diag}(w_e)B^T$
i dowolnego jej minora głównego $L^{(r)}$
uzyskanego przez skreślenie wiersza i kolumny $r$
zachodzi
\[
\det L^{(r)}=\sum_{T\text{ drzewo rozpinające }G}
\ \prod_{e\in E(T)}w_e.
\]
Tożsamość obowiązuje także dla wag ujemnych.
\end{stwierdzenie}
\begin{proof}
Niech $B_r$ oznacza macierz powstałą z $B$ przez
usunięcie wiersza $r$. Z formuły Cauchy'ego--Bineta
mamy
\[
\det(B_r\operatorname{diag}(w_e)B_r^T)
=\sum_{F\subseteq E,\ |F|=n-1}
\det(B_{r,F})^2\prod_{e\in F}w_e.
\]
Tutaj $B_{r,F}$ jest kwadratową podmacierzą $B_r$
o kolumnach indeksowanych przez $F$. Drugi czynnik
we wzorze Cauchy'ego--Bineta ma minor równy
$\det(B_{r,F}^T)\prod_{e\in F}w_e$.

Jeżeli $(V,F)$ nie jest spójny, suma wierszy
$B_{r,F}$ należących do składowej niezawierającej $r$
jest zerowa: obie incydencje krawędzi wewnętrznej
znoszą się, a krawędzi wychodzących ze składowej
nie ma. Jest to nietrywialna zależność liniowa.
Jeżeli $(V,F)$ jest spójny, zawiera drzewo rozpinające:
usuwamy pętle, utożsamiamy krawędzie równoległe,
wybieramy drzewo grafu prostego i podnosimy jego
krawędzie do $F$. Ponieważ $|F|=n-1$, drzewo
wykorzystuje wszystkie krawędzie $F$.
Dla $n>1$ wybieramy liść różny od $r$ i rozwijamy
wyznacznik względem jego wiersza. Jedyny niezerowy
element tego wiersza jest równy $\pm1$; pozostaje
macierz drzewa z usuniętym liściem i jego krawędzią.
Indukcja prowadzi do pustej macierzy o wyznaczniku $1$.
Stąd $\det B_{r,F}=\pm1$. Dla $n=1$ jedyne drzewo
ma pusty zbiór krawędzi i wagę $1$.
Otrzymujemy rozszerzenie twierdzenia
\ref{thm:tg-kirchhoff} na multigrafy i dowolne wagi,
w tym zerowe. Wynik nie zależy od wyboru $r$.
\end{proof}

\begin{stwierdzenie}[Drzewa rozpinające planarnych grafów dualnych]
\label{thm:tg-tait-dual-trees}
Niech $G$ będzie skończonym, spójnym grafem
zanurzonym komórkowo w $S^2$, z dopuszczeniem
pętli i krawędzi wielokrotnych, a $G^*$
jego grafem dualnym. Wtedy
\[
\tau(G)=\tau(G^*).
\]
Dokładniej, jeśli $T\subset E(G)$ jest drzewem
rozpinającym, to krawędzie dualne do
$E(G)\setminus E(T)$ tworzą drzewo rozpinające
grafu $G^*$.
\end{stwierdzenie}

\begin{proof}
Twierdzenie \ref{thm:tg-dual-spanning} daje drzewo
rozpinające $T^*=\{e^*:e\notin E(T)\}$ grafu dualnego,
także gdy grafy zawierają pętle i krawędzie równoległe.
Stosujemy to samo twierdzenie do $G^*$ i $T^*$.
Z twierdzenia \ref{thm:tg-cellular-dual} oraz stwierdzenia
\ref{thm:tg-dual-basic} identyfikujemy podwójny graf
dualny z grafem pierwotnym, zachowując odpowiedniość
krawędzi. Wobec tego dopełnienie zbioru krawędzi $T^*$
po ponownej dualizacji jest dokładnie zbiorem $E(T)$.
Otrzymujemy wzajemnie odwrotne odwzorowania między
zbiorami drzew rozpinających $G$ i $G^*$.
Są to zbiory skończone, więc ich liczności są równe.
\end{proof}

\begin{wniosek}[Ważona dualność drzew rozpinających]
\label{cor:tg-tait-weighted-dual}
Niech $G$ będzie spójnym grafem komórkowo zanurzonym
w $S^2$, a $G^*$ jego grafem dualnym.
Niech każdej krawędzi $e\in E(G)$
przypisano niezerową wagę $w_e\in\mathbb R$.
Krawędzi dualnej $e^*$ przypiszmy wagę
$w^*_{e^*}=w_e^{-1}$.
Oznaczmy przez
$\tau_w(G)=\sum_{T}\prod_{e\in E(T)}w_e$
ważoną sumę po drzewach rozpinających.
Wówczas
\[
\tau_w(G)=\Bigl(\prod_{e\in E(G)}w_e\Bigr)
\tau_{w^*}(G^*).
\]
Tożsamość jest prawdziwa również dla wag ujemnych.
\end{wniosek}

\begin{proof}
Ze stwierdzenia~\ref{thm:tg-tait-dual-trees}
drzewu rozpinającemu $T\subset G$
odpowiada drzewo dualne
$T^*=\{e^*:e\notin T\}$.
Dla każdej takiej pary mamy
\[
\prod_{e\in T}w_e
=\Bigl(\prod_{e\in E(G)}w_e\Bigr)
\prod_{e^*\in T^*}w^*_{e^*}.
\]
Sumując po wszystkich drzewach $T$ i korzystając
z bijekcji $T\mapsto T^*$, otrzymujemy tezę.
Wymóg $w_e\ne0$ służy wyłącznie
zdefiniowaniu wag odwrotnych; dla wag zerowych
pozostaje wcześniejsza, bezpośrednia
formuła z minorów podpisanego laplasjanu.
\end{proof}

\begin{uwaga}[Wagi ujemne nie zliczają drzew]
Tożsamość z poprzedniego stwierdzenia
nie pozwala usuwać wag w formule
Kirchhoffa. Na przykład graf o dwóch
wierzchołkach i dwóch krawędziach
równoległych, z wagami $+1$ i $-1$,
ma dwa drzewa rozpinające, ale jego
podpisany laplasjan jest macierzą zerową.
Suma ich wag wynosi $1+(-1)=0$.
Jednakowe znaki wag, jak w przypadku diagramu alternującego,
są warunkiem wystarczającym równości
$|\det\mathcal G_D^{(r)}|=\tau(G_B(D))$.
Dokładniej, równość zachodzi wtedy i tylko wtedy, gdy wszystkie
iloczyny wag po drzewach rozpinających mają ten sam znak:
jeżeli $p$ z nich wynosi $+1$, a $q$ wynosi $-1$, to
$|p-q|=p+q$ wtedy i tylko wtedy, gdy $pq=0$.
Nie wymaga to jednakowych znaków wszystkich krawędzi;
na przykład w drzewie o dowolnych wagach $\pm1$ istnieje
tylko jedno drzewo rozpinające.
\end{uwaga}

\begin{stwierdzenie}[Specjalizacja nawiasu a podpisana suma drzew]
\label{prop:tg-tait-bracket-det}
Niech $D$ będzie spójnym diagramem splotu z co najmniej
jednym skrzyżowaniem, niekoniecznie alternującym.
Niech $G=G_B(D)$, $v=|V(G)|$, a $\eta(e)\in\{-1,1\}$
będzie znakiem z definicji podpisanego laplasjanu.
Dla $\zeta=\exp(\mathrm i\pi/4)$ zachodzi
\[
 \langle D\rangle\big|_{A=\zeta}
 =\zeta^{-\sum_{e\in E(G)}\eta(e)}
 \mathrm i^{v-1}\det\mathcal G_D^{(r)}.
\]
W szczególności moduł lewej strony jest równy
$|\det\mathcal G_D^{(r)}|$.
\end{stwierdzenie}

\begin{proof}
Każdy stan wygładzeń kodujemy zbiorem $F$ krawędzi,
dla których wygładzenie łączy czarne sektory.
Dla $\eta(e)=1$ jest to wygładzenie A, a dla
$\eta(e)=-1$ wygładzenie B. Dla krawędzi spoza $F$
typ wygładzenia jest przeciwny. Współczynnik stanu
jest więc równy
\[
 A^{-\sum_e\eta(e)+2\sum_{e\in F}\eta(e)}.
\]
Argument z dyskami i pasami z dowodu twierdzenia
\ref{thm:tg-tait-tutte} nie wymaga alternującego
diagramu, jeśli $F$ oznacza właśnie krawędzie
łączące czarne sektory. Liczba okręgów stanu
wynosi $b(F)=2k(F)-v+|F|=k(F)+\nu_G(F)$,
gdzie $\nu_G(F)=|F|-v+k(F)\geq0$.
Ponieważ $k(F)\geq1$, równość $b(F)=1$ zachodzi
dokładnie wtedy, gdy $k(F)=1$ i $\nu_G(F)=0$,
czyli gdy $(V(G),F)$ jest drzewem rozpinającym
(por. argument dla spójnego podgrafu o $v-1$
krawędziach w dowodzie \ref{prop:tg-tait-signed-tree}).

Dla $A=\zeta$ mamy $\zeta^2=\mathrm i$,
$\zeta^{-2}=-\mathrm i$ i $\delta=0$.
Czynnik $\delta^{b(F)-1}$ znika dla $b(F)>1$
i jest równy $1$ dla $b(F)=1$ (wykładnik zero).
Ponadto $\mathrm i^{\eta(e)}=\mathrm i\eta(e)$
dla obu możliwych znaków. Każde drzewo ma $v-1$
krawędzi, więc po podstawieniu pozostaje
\[
 \langle D\rangle\big|_{A=\zeta}
 =\zeta^{-\sum_e\eta(e)}\mathrm i^{v-1}
   \sum_{T\text{ drzewo rozpinające }G}
   \prod_{e\in E(T)}\eta(e).
\]
Stwierdzenie \ref{prop:tg-tait-signed-tree}
daje wskazany wyznacznik. Czynnik przed sumą
ma moduł jeden. Jest to niezależny rachunek
kombinatoryczny, bez korzystania z prezentacji
homologii nakrycia rozgałęzionego.
\end{proof}

\begin{twierdzenie}[Goeritz; identyfikacja wyznacznika]
Niech $K\subset S^3$ będzie węzłem, a $D$ jego
spójnym regularnym diagramem z co najmniej jednym
skrzyżowaniem. Niech $\mathcal G_D^{(r)}$ będzie dowolną zredukowaną
macierzą Goeritza zdefiniowaną powyżej. Wtedy
\[
\det(K)=|\Delta_K(-1)|=|\det\mathcal G_D^{(r)}|.
\]
W szczególności, dla alternującego diagramu węzła,
po wyborze kolorowania, przy którym wszystkie wagi
są jednakowe, otrzymujemy
\[
\det(K)=\tau(G_B(D))=\tau(G_W(D)),
\]
gdzie $\tau(G)$ oznacza liczbę drzew rozpinających
grafu $G$. Ostatnią równość uzasadnia
Stwierdzenie~\ref{thm:tg-tait-dual-trees}.
\end{twierdzenie}
\begin{proof}[Dowód z klasycznej prezentacji Goeritza]
Korzystamy z zewnętrznego twierdzenia o prezentacji
Goeritza: zredukowana macierz Goeritza spójnego
diagramu prezentuje $H_1(\Sigma_2(K);\mathbb Z)$.
Sprawdźmy zgodność konwencji z rysunkiem 1 i wzorem
(2.1) w pracy Silvera, Traldiego i Williams (cytowanej
niżej). Ich obszary niezacieniowane odpowiadają naszym
czarnym. Na lewym rysunku pasmo górne biegnie od lewego
górnego do prawego dolnego końca, a obszary niezacieniowane
zajmują sektory górny i dolny. Ich liczba incydencji
wynosi $+1$, zgodnie z naszym wygładzeniem A. Na prawym
rysunku przy tym samym kolorowaniu zamieniono pasma
górne i dolne; liczba incydencji wynosi $-1$, również
zgodnie z naszą definicją. Obrót o ćwierć obrotu
sprowadza pozostałe dwa przypadki do tych już sprawdzonych.
Zgodne są także wpisy pozadiagonalne i diagonalne,
w tym brak wkładu pętli. Wybierając punkt w czarnym
obszarze odpowiadającym $r$ jako punkt w nieskończoności,
otrzymujemy dokładnie ich redukcję względem obszaru
nieograniczonego. Oba grafy Taita są spójne, więc
nie występuje dodatkowa część wolna pochodząca
z niespójności grafu szachownicowego. Zatem $H_1(\Sigma_2(K);\mathbb Z)\cong
\operatorname{coker}\mathcal G_D^{(r)}$.
Z sekcji \ref{sec:kt-branched} grupa ta jest skończona
i ma rząd $|\Delta_K(-1)|$. Postać normalna Smitha
daje równość rzędu kokernu z modułem wyznacznika.
Dla diagramu alternującego lemat
\ref{lem:tg-tait-alternating} pozwala wybrać wszystkie
wagi równe $+1$. Stwierdzenie
\ref{prop:tg-tait-signed-tree} daje liczbę drzew
grafu czarnego, a \ref{thm:tg-tait-dual-trees}
liczbę drzew grafu białego.
\end{proof}

\begin{uwaga}[Zakres dowodu]
Zewnętrznym składnikiem jest prezentacja Goeritza,
nie wzór sumy po drzewach. Zob. W. B. R. Lickorish,
\emph{An Introduction to Knot Theory}, Springer, 1997,
rozdział 9, \emph{Cyclic Branched Covers and the Goeritz Matrix},
oraz C. McA. Gordon, R. A. Litherland,
\emph{On the signature of a link}, Invent. Math. 47 (1978),
53--69; D. S. Silver, L. Traldi, S. G. Williams,
\emph{Goeritz and Seifert matrices from Dehn presentations},
Osaka J. Math. 57 (2020), 663--677, w szczególności
s. 663--664 (rysunek 1 i wzór (2.1)) oraz twierdzenie 4.1.
Nie wyprowadzamy tutaj topologicznej konstrukcji
prezentacji. Dla diagramu bez skrzyżowań mamy węzeł
trywialny i $\Delta_U=1$; zredukowana macierz
jednowierzchołkowego grafu jest pusta i ma wyznacznik $1$.
\end{uwaga}

\begin{przyklad}[Wyznacznik trójlistnika bez macierzy Seiferta]
Dla $G=C_3$ wszystkie wagi można przyjąć za $+1$.
Laplasjan grafu ma macierz
\[
\mathcal G_D=\begin{pmatrix}
 2&-1&-1\\ -1&2&-1\\ -1&-1&2
\end{pmatrix}.
\]
Po skreśleniu trzeciego wiersza i kolumny dostajemy
$\mathcal G_D^{(3)}=\begin{psmallmatrix}2&-1\\-1&2\end{psmallmatrix}$,
więc $\det(K)=|4-1|=3$.
Ta sama liczba jest liczbą trzech drzew rozpinających
trójkątnego grafu Taita i zgadza się z wyznacznikiem
wielomianu Alexandera $t^2-t+1$ przy $t=-1$.
\end{przyklad}

\begin{przyklad}[Podpisana macierz Goeritza po zmianie skrzyżowania]
\label{ex:tg-goeritz-crossing-change}
Rozważmy graf Taita standardowego
alternującego diagramu trójlistnika:
jest on cyklem $C_3$ o trzech krawędziach.
Dla wybranego kolorowania wagi
wszystkich krawędzi są równe $+1$.
Zmiana przejścia górą na przejście dołem
w jednym skrzyżowaniu zamienia w tym
miejscu wygładzenia A i B, a więc
zmienia znak jednej wagi.
Dla wag $(1,1,-1)$ otrzymujemy
\[
\mathcal G=
\begin{pmatrix}
0&-1&1\\
-1&2&-1\\
1&-1&0
\end{pmatrix},
\qquad
\det\mathcal G^{(3)}
=\det\begin{pmatrix}0&-1\\-1&2\end{pmatrix}
=-1.
\]
Równoważnie sumowanie po trzech drzewach
rozpinających daje
$\tau_w(C_3)=1-1-1=-1$.
Po zastosowaniu twierdzenia Goeritza
wyznacznik zmienionego węzła wynosi więc $1$,
mimo że \emph{niepodpisany} graf Taita
wciąż ma trzy drzewa rozpinające.

W standardowym diagramie trójlistnika jako domknięcia
dwupasmowego warkocza $\sigma_1^3$ zmiana jednego
skrzyżowania daje słowo z dwoma czynnikami $\sigma_1$
i jednym $\sigma_1^{-1}$. Skracając sąsiednią parę
odwrotnych skrzyżowań ruchem Reidemeistera II,
otrzymujemy domknięcie $\sigma_1$, czyli diagram
okręgu z jednym skrętem, usuwanym ruchem I. Zbieżność
$\det(U)=1$ z powyższym rachunkiem
jest kontrolą zgodności konwencji znaków,
ale nie dowodzi twierdzenia Goeritza
dla wszystkich diagramów.
\end{przyklad}

\begin{uwaga}[Literatura]
Związek wielomianu Jonesa i wielomianu Tutte'a dla diagramów
alternujących jest klasycznym wynikiem Thistlethwaite'a.
Metodę wyznaczników macierzy Goeritza i jej związek z
nakryciami rozgałęzionymi opisują podręczniki
W. B. R. Lickorisha i D. Rolfsena.
\end{uwaga}


